\begin{longtable}[c]{|p{2cm}<{\centering}|p{2cm}<{\centering}|p{2cm}<{\centering}|p{2cm}<{\centering}|p{2cm}<{\centering}|p{4cm}|}
    \caption{SEG.4 - SEG.0 的位信息}
    \label{tab:twai-bit-info-seg4-seg0}\\
    \hline
    \rowcolor{lightgray}
   \textbf{Bit SEG.4} & \textbf{Bit SEG.3} & \textbf{Bit SEG.2} & \textbf{Bit SEG.1} & \textbf{Bit SEG.0} & \textbf{描述} \\\hline
   \endfirsthead
   
   \multicolumn{6}{c}{\bfseries \tablename\ \thetable{} -- 接上页}\\\hline
   
   \rowcolor{lightgray}
   \textbf{Bit SEG.4} & \textbf{Bit SEG.3} & \textbf{Bit SEG.2} & \textbf{Bit SEG.1} & \textbf{Bit SEG.0} & \textbf{描述} \\\hline
   \endhead
   \multicolumn{6}{r}{\textbf{见下页}} \\
   \endfoot
   \endlastfoot
   
    0 & 0 & 0 & 1 & 1 & 帧起始
    \\\hline
    0 & 0 & 0 & 1 & 0 & ID.28 \verb+~+ ID.21
    \\\hline
    0 & 0 & 1 & 1 & 0 & ID.20 \verb+~+ ID.18
    \\\hline
    0 & 0 & 1 & 0 & 0 & bit SRTR
    \\\hline
    0 & 0 & 1 & 0 & 1 & bit IDE
    \\\hline
    0 & 0 & 1 & 1 & 1 & ID.17 \verb+~+ ID.13
    \\\hline
    0 & 1 & 1 & 1 & 1 & ID.12 \verb+~+ ID.5
    \\\hline
    0 & 1 & 1 & 1 & 0 & ID.4 \verb+~+ ID.0
    \\\hline
    0 & 1 & 1 & 0 & 0 & bit RTR
    \\\hline
    0 & 1 & 1 & 0 & 1 & 保留位 1
    \\\hline
    0 & 1 & 0 & 0 & 1 & 保留位 0
    \\\hline
    0 & 1 & 0 & 1 & 1 & 数据长度代码
    \\\hline
    0 & 1 & 0 & 1 & 0 & 数据域
    \\\hline
    0 & 1 & 0 & 0 & 0 & CRC 序列
    \\\hline
    1 & 1 & 0 & 0 & 0 & CRC 分界符
    \\\hline
    1 & 1 & 0 & 0 & 1 & 确认槽
    \\\hline
    1 & 1 & 0 & 1 & 1 & 确认分界符
    \\\hline
    1 & 1 & 0 & 1 & 0 & 帧结束
    \\\hline
    1 & 0 & 0 & 1 & 0 & 间歇域
    \\\hline
    1 & 0 & 0 & 0 & 1 & 主动错误标志
    \\\hline
    1 & 0 & 1 & 1 & 0 & 被动错误标志
    \\\hline
    1 & 0 & 0 & 1 & 1 & 兼容显性位
    \\\hline
    1 & 0 & 1 & 1 & 1 & 错误分界符
    \\\hline
    1 & 1 & 1 & 0 & 0 & 过载标志
    \\\hline
\end{longtable}
\vspace{-1.5em}
{\bfseries 说明：}\\
\begin{itemize}
    \item Bit SRTR: 标准格式 RTR bit。
    \item Bit IDE: 标识符扩展位。0 表示标准格式。
\end{itemize}